%!TEX program = lualatex
% Arquivo: relatorio_pratica_4_NaOH_HCl.tex

\documentclass[a4paper,12pt]{article}

% =========================================================
% Idioma e tipografia (LuaLaTeX)
% =========================================================
\usepackage[brazilian]{babel}
\usepackage{fontspec}
\setmainfont{Latin Modern Roman}
\usepackage{microtype}

% =========================================================
% Layout
% =========================================================
\usepackage[top=3cm,bottom=2.5cm,left=3cm,right=2.5cm]{geometry}
\usepackage{setspace}
\onehalfspacing
\usepackage{indentfirst}
\setlength{\parindent}{1.25cm}

% =========================================================
% Matemática, unidades e fórmulas químicas
% =========================================================
\usepackage{amsmath,amssymb}
\usepackage{siunitx}
\sisetup{
  locale = DE,
  output-decimal-marker = {,},
  per-mode = symbol,
  exponent-product = \cdot
}
\usepackage[version=4]{mhchem}

% =========================================================
% Tabelas, listas, figuras e desenhos
% =========================================================
\usepackage{booktabs}
\usepackage{array}
\usepackage{enumitem}
\setlist[itemize]{noitemsep,topsep=3pt}
\setlist[enumerate]{noitemsep,topsep=3pt}
\usepackage{graphicx}
\usepackage{float}
\graphicspath{{./}{./imagens/}{./figuras/}}

\usepackage{xcolor}
\usepackage{tikz}
\usetikzlibrary{arrows.meta,calc,positioning,shapes.geometric}
\usepackage[most]{tcolorbox}

% =========================================================
% Links
% =========================================================
\usepackage{hyperref}
\usepackage{xurl}
\hypersetup{
  colorlinks=true,
  linkcolor=black,
  urlcolor=blue,
  citecolor=black,
  pdfauthor={Israel Carneiro; José Mauro Xavier Elsas; Miguel Gayani},
  pdftitle={Prática 4 -- Preparo e padronização de solução de NaOH}
}

% =========================================================
% Cores e comandos
% =========================================================
\definecolor{quadroazul}{RGB}{40,91,140}
\definecolor{quadrovermelho}{RGB}{177,69,60}
\definecolor{quadrocinza}{RGB}{245,245,245}

\newcommand{\subtitulo}[1]{%
  \par\medskip
  {\color{quadroazul}\bfseries\large #1}\par\smallskip
}

% =========================================================
% Dados da capa
% =========================================================
\newcommand{\Instituicao}{Universidade do Estado do Rio de Janeiro -- UERJ}
\newcommand{\Curso}{Química VIII / Oceanografia}
\newcommand{\Professor}{Professor: Ítalo Vieira}
\newcommand{\Aluno}{Alunos: Israel Carneiro\\
José Mauro Xavier Elsas\\
Miguel Gayani}
\newcommand{\Cidade}{Rio de Janeiro}
\newcommand{\Ano}{2026}

% =========================================================
% Dados numéricos da prática
% =========================================================
% A concentração abaixo foi adotada como 0,1000 mol/L por não constar,
% no material disponível, o valor final da padronização da Prática 3.
% Quando o valor padronizado estiver disponível, substituir aqui e
% recalcular os resultados dependentes.
\newcommand{\ConcHCl}{0.1000}

\begin{document}

% =========================================================
% Capa
% =========================================================
\begin{titlepage}
\centering
{\large \Instituicao \par}
\vspace{0.5cm}
{\large \Curso \par}
\vspace{1.0cm}

{\LARGE \textbf{Prática 4} \par}
\vspace{0.6cm}
{\Large \textbf{PREPARO E PADRONIZAÇÃO DE SOLUÇÃO\\
DE \ce{NaOH} \qty{0.1}{\mole\per\litre}} \par}

\vspace{1cm}

\IfFileExists{uerj.jpg}{%
  \begin{figure}[H]
    \centering
    \includegraphics[width=0.3\linewidth]{uerj.jpg}
  \end{figure}
}{%
  \begin{tcolorbox}[
    width=0.32\textwidth,
    colback=white,
    colframe=black!35,
    boxrule=0.5pt,
    arc=0mm
  ]
  \centering
  \vspace{1.2cm}
  \textit{Logotipo da UERJ}
  \vspace{1.2cm}
  \end{tcolorbox}
}

\vfill

\begin{flushright}
\Aluno\\
\Professor
\end{flushright}

\vspace{1.0cm}
{\Cidade{} -- \Ano \par}
\end{titlepage}

% =========================================================
\section{Introdução}

O hidróxido de sódio é uma das bases mais empregadas em análise volumétrica.
Entretanto, o \ce{NaOH} não é considerado um padrão primário, porque é
higroscópico e absorve dióxido de carbono da atmosfera. A reação com o
\ce{CO2} produz carbonato e altera gradualmente a composição da solução:

\begin{equation}
  \ce{CO2(g) + 2OH-(aq) -> CO3^2-(aq) + H2O(l)}.
  \label{eq:co2}
\end{equation}

Por essa razão, uma solução de \ce{NaOH} preparada por diluição apresenta
concentração apenas aproximada e deve ser padronizada antes de ser usada em
determinações quantitativas. Nesta prática, a solução básica foi padronizada
por titulação com a solução de \ce{HCl} preparada na Prática 3.

A reação de neutralização entre o ácido clorídrico e o hidróxido de sódio é:

\begin{equation}
  \ce{HCl(aq) + NaOH(aq) -> NaCl(aq) + H2O(l)}.
  \label{eq:neutralizacao}
\end{equation}

Como os coeficientes estequiométricos de \ce{HCl} e \ce{NaOH} são iguais a
um, a relação entre as quantidades de matéria no ponto de equivalência é de
$1:1$.

% =========================================================
\section{Objetivos}

\begin{itemize}
  \item preparar aproximadamente \qty{250}{\milli\litre} de solução de
        \ce{NaOH} \qty{0.1}{\mole\per\litre} por diluição de uma solução
        concentrada;
  \item padronizar a solução preparada por titulação com a solução de
        \ce{HCl} obtida na prática anterior;
  \item determinar a concentração molar experimental da solução de
        \ce{NaOH};
  \item distinguir o volume da alíquota de \ce{NaOH} do volume total
        existente no Erlenmeyer após a adição de água destilada.
\end{itemize}

% =========================================================
\section{Material utilizado}

\begin{enumerate}[label=\alph*.]
  \item béquer;
  \item cilindro graduado de \qty{500}{\milli\litre};
  \item pipeta graduada para medir aproximadamente \qty{2.0}{\milli\litre};
  \item pipeta de \qty{25.00}{\milli\litre};
  \item bureta de \qty{50.00}{\milli\litre};
  \item suporte universal e garra para bureta;
  \item três Erlenmeyers de \qty{250}{\milli\litre};
  \item pisseta com água destilada;
  \item frasco para armazenamento da solução preparada.
\end{enumerate}

% =========================================================
\section{Reagentes e soluções}

\begin{enumerate}[label=\alph*.]
  \item solução concentrada de \ce{NaOH}, aproximadamente
        \qty{12.4}{\mole\per\litre};
  \item solução de \ce{HCl} preparada na Prática 3, adotada neste relatório
        como \qty{\ConcHCl}{\mole\per\litre};
  \item água destilada;
  \item solução indicadora de fenolftaleína.
\end{enumerate}

\begin{tcolorbox}[
  title=Observação sobre a concentração do \ce{HCl},
  colback=quadrocinza,
  colframe=black!65,
  boxrule=0.6pt,
  arc=1mm
]
O valor numérico empregado nos cálculos foi
\qty{\ConcHCl}{\mole\per\litre}. Caso a concentração padronizada obtida na
Prática 3 seja diferente, esse valor deve substituir a concentração nominal,
e a concentração calculada de \ce{NaOH} deve ser atualizada.
\end{tcolorbox}

% =========================================================
\section{Procedimento experimental}

\subsection{Preparo da solução aproximada de \ce{NaOH}}

Pretendia-se preparar \qty{250}{\milli\litre} de solução de \ce{NaOH} com
concentração aproximada de \qty{0.1}{\mole\per\litre}, a partir de uma solução
concentrada de \qty{12.4}{\mole\per\litre}. O volume necessário da solução
concentrada foi calculado pela relação de diluição:

\begin{equation}
  C_1V_1=C_2V_2.
  \label{eq:diluicao}
\end{equation}

Isolando $V_1$:

\begin{align}
  V_1
    &=\frac{C_2V_2}{C_1} \nonumber\\
    &=\frac{
       \left(\qty{0.1}{\mole\per\litre}\right)
       \left(\qty{0.250}{\litre}\right)}
       {\qty{12.4}{\mole\per\litre}} \nonumber\\
    &=\qty{0.002016}{\litre} \nonumber\\
    &=\qty{2.016}{\milli\litre}
    \approx \boxed{\qty{2.0}{\milli\litre}}.
  \label{eq:volume_naoh_conc}
\end{align}

Foram colocados aproximadamente \qty{100}{\milli\litre} de água destilada no
cilindro graduado. Em seguida, adicionaram-se aproximadamente
\qty{2.0}{\milli\litre} da solução concentrada de \ce{NaOH}. Após a mistura e
o resfriamento, o volume foi completado até \qty{250}{\milli\litre} com água
destilada, e a solução foi homogeneizada.

\begin{figure}[H]
\centering
\begin{tikzpicture}[>=Latex,thick,font=\small]
  % pipeta
  \begin{scope}[shift={(-4.2,2.3)},rotate=-18]
    \draw[quadrovermelho] (-1.6,0) -- (-0.45,0);
    \draw[quadrovermelho]
      (-0.45,0) .. controls (-0.25,0.35) and (0.25,0.35) .. (0.45,0)
      .. controls (0.25,-0.35) and (-0.25,-0.35) .. (-0.45,0);
    \draw[quadrovermelho] (0.45,0) -- (1.8,0);
    \node[quadrovermelho,above=4mm,rotate=18] at (0,0.25)
      {\qty{2.0}{\milli\litre} de \ce{NaOH} concentrado};
  \end{scope}

  \draw[->,quadrovermelho] (-2.4,1.8) -- (-0.9,1.25);

  % cilindro
  \draw[quadrovermelho,fill=white] (-0.7,0) rectangle (0.3,3.4);
  \draw[quadrovermelho] (-0.85,3.4) -- (0.45,3.4);
  \draw[quadrovermelho] (-0.9,0) -- (0.5,0);
  \fill[quadroazul!16] (-0.67,0.05) rectangle (0.27,1.30);
  \draw[quadroazul] (-0.67,1.30) -- (0.27,1.30);
  \foreach \y in {0.4,0.8,1.2,1.6,2.0,2.4,2.8,3.2}
    \draw[quadrovermelho] (-0.7,\y) -- (-0.45,\y);
  \node[below=2mm,align=center] at (-0.2,0)
    {cilindro graduado\\de \qty{500}{\milli\litre}};
  \node[quadroazul,right=5mm,align=left] at (0.3,1.3)
    {aproximadamente\\\qty{100}{\milli\litre} de água};

  \draw[->,black!65] (1.2,1.6) -- (3.0,1.6);

  % recipiente final
  \draw[quadrovermelho,fill=white] (3.3,0) rectangle (4.3,3.4);
  \draw[quadrovermelho] (3.15,3.4) -- (4.45,3.4);
  \draw[quadrovermelho] (3.1,0) -- (4.5,0);
  \fill[quadroazul!16] (3.33,0.05) rectangle (4.27,2.1);
  \draw[quadroazul] (3.33,2.1) -- (4.27,2.1);
  \node[below=2mm,align=center] at (3.8,0)
    {completar até\\\qty{250}{\milli\litre}};
\end{tikzpicture}
\caption{Esquema do preparo da solução aproximada de \ce{NaOH}.}
\label{fig:preparo_naoh}
\end{figure}

\subsection{Padronização da solução de \ce{NaOH} com \ce{HCl}}

A bureta foi rinsada e preenchida com a solução de \ce{HCl}. Para cada
determinação, transferiu-se uma alíquota de \qty{25.00}{\milli\litre} da
solução de \ce{NaOH} para um Erlenmeyer de \qty{250}{\milli\litre}.
Adicionaram-se aproximadamente \qty{25}{\milli\litre} de água destilada e
duas gotas de fenolftaleína.

A solução no Erlenmeyer apresentou coloração rosa em meio básico. O
\ce{HCl} foi adicionado lentamente, sob agitação, até o desaparecimento
persistente da coloração rosa, adotado como ponto final da titulação. O
procedimento foi realizado em triplicata.

\begin{figure}[H]
\centering
\begin{tikzpicture}[>=Latex,thick,font=\small]
  % bureta
  \draw[quadroazul] (-0.15,1.8) rectangle (0.15,5.5);
  \draw[quadroazul] (-0.30,5.5) -- (0.30,5.5);
  \draw[quadroazul] (-0.15,1.8) -- (0,1.3) -- (0.15,1.8);
  \draw[quadroazul] (0,1.3) -- (0,0.9);
  \node[quadroazul,right=5mm] at (0.15,4.7)
    {\ce{HCl} $\approx$ \qty{0.1}{\mole\per\litre}};

  % gota
  \fill[quadroazul!50] (0,0.65) circle (0.07);

  % erlenmeyer
  \draw[quadrovermelho]
    (-0.35,0.35) .. controls (-1.25,-0.15) and (-1.15,-1.05) .. (-0.55,-1.35)
    -- (0.55,-1.35)
    .. controls (1.15,-1.05) and (1.25,-0.15) .. (0.35,0.35)
    -- (0.35,1.10) -- (-0.35,1.10) -- cycle;
  \fill[quadrovermelho!12]
    (-0.70,-0.95) .. controls (-0.90,-0.65) and (-0.85,-0.35) .. (-0.55,-0.20)
    -- (0.55,-0.20)
    .. controls (0.85,-0.35) and (0.90,-0.65) .. (0.70,-0.95) -- cycle;
  \node[align=center,below=5mm] at (0,-1.35)
    {\qty{25.00}{\milli\litre} de \ce{NaOH}\\
     + \qty{25}{\milli\litre} de água\\
     + 2 gotas de fenolftaleína};

  \draw[->,quadrovermelho] (2.2,0.6) -- (0.55,-0.1);
  \node[quadrovermelho,align=left,right] at (2.25,0.6)
    {rosa $\longrightarrow$ incolor\\no ponto final};
\end{tikzpicture}
\caption{Esquema da padronização da solução de \ce{NaOH} com \ce{HCl}.}
\label{fig:titulacao_naoh}
\end{figure}

% =========================================================
\section{Resultados e cálculos}

\subsection{Volumes de \ce{HCl} consumidos}

Os volumes registrados nas três determinações foram:

\begin{table}[H]
  \centering
  \caption{Volumes de solução de \ce{HCl} consumidos nas titulações.}
  \label{tab:volumes}
  \begin{tabular}{
    c
    S[table-format=2.2]
  }
    \toprule
    \textbf{Determinação}
      & {\textbf{Volume de \ce{HCl}} (\si{\milli\litre})} \\
    \midrule
    1 & 30,70 \\
    2 & 30,30 \\
    3 & 30,20 \\
    \bottomrule
  \end{tabular}
\end{table}

O volume médio foi:

\begin{align}
  \overline{V}_{\ce{HCl}}
    &=\frac{V_1+V_2+V_3}{3} \nonumber\\
    &=\frac{30.70+30.30+30.20}{3}\,\si{\milli\litre} \nonumber\\
    &=\boxed{\qty{30.40}{\milli\litre}}.
  \label{eq:volume_medio}
\end{align}

O desvio-padrão amostral dos volumes foi de aproximadamente
\qty{0.26}{\milli\litre}, e o desvio-padrão relativo foi de
aproximadamente \qty{0.87}{\percent}. Esses valores indicam boa
proximidade entre as três determinações.

\subsection{Por que os \qty{25}{\milli\litre} de água não entram no cálculo}

A alíquota retirada da solução original continha
\qty{25.00}{\milli\litre} de solução de \ce{NaOH}. A quantidade de matéria
de base presente nessa alíquota era:

\begin{equation}
  n_{\ce{NaOH}}
  =C_{\ce{NaOH}}V_{\ce{NaOH}}.
  \label{eq:mols_naoh}
\end{equation}

Após a adição de aproximadamente \qty{25}{\milli\litre} de água
destilada, o volume total no Erlenmeyer passou a ser aproximadamente
\qty{50}{\milli\litre}. Entretanto, a água não adicionou \ce{NaOH} e,
portanto, não alterou a quantidade de matéria da base:

\begin{equation}
  n_{\ce{NaOH},\text{antes}}
  =
  n_{\ce{NaOH},\text{depois}}.
  \label{eq:mols_constantes}
\end{equation}

Considerando a diluição ideal:

\begin{equation}
  C_{\text{diluída}}
  =
  C_{\ce{NaOH}}
  \frac{\qty{25.00}{\milli\litre}}
       {\qty{50.00}{\milli\litre}}
  =
  \frac{C_{\ce{NaOH}}}{2}.
  \label{eq:concentracao_diluida}
\end{equation}

Assim:

\begin{align}
  n_{\ce{NaOH}}
    &=C_{\text{diluída}}V_{\text{total}} \nonumber\\
    &=\frac{C_{\ce{NaOH}}}{2}
      \left(\qty{50.00}{\milli\litre}\right) \nonumber\\
    &=C_{\ce{NaOH}}
      \left(\qty{25.00}{\milli\litre}\right).
  \label{eq:demonstracao_diluicao}
\end{align}

Portanto, os \qty{25}{\milli\litre} de água contam para o volume total
presente no Erlenmeyer, mas não são somados ao volume da alíquota de
\ce{NaOH} na relação estequiométrica. No cálculo da concentração da
solução original, utiliza-se:

\begin{equation}
  \boxed{V_{\ce{NaOH}}=\qty{25.00}{\milli\litre}}.
  \label{eq:volume_correto}
\end{equation}

A água foi adicionada para aumentar o volume no Erlenmeyer, facilitar a
agitação, lavar as paredes internas e tornar mais nítida a observação da
mudança de cor da fenolftaleína.

\subsection{Cálculo da concentração da solução de \ce{NaOH}}

No ponto de equivalência, devido à proporção estequiométrica $1:1$ da
Equação~\ref{eq:neutralizacao}:

\begin{equation}
  n_{\ce{HCl}}=n_{\ce{NaOH}}.
\end{equation}

Consequentemente:

\begin{equation}
  C_{\ce{HCl}}V_{\ce{HCl}}
  =
  C_{\ce{NaOH}}V_{\ce{NaOH}}.
  \label{eq:relacao_titulacao}
\end{equation}

Isolando a concentração de \ce{NaOH}:

\begin{equation}
  C_{\ce{NaOH}}
  =
  \frac{C_{\ce{HCl}}V_{\ce{HCl}}}
       {V_{\ce{NaOH}}}.
  \label{eq:cnaoh}
\end{equation}

Usando o volume médio de \ce{HCl}, a alíquota de
\qty{25.00}{\milli\litre} de \ce{NaOH} e
$C_{\ce{HCl}}=\qty{\ConcHCl}{\mole\per\litre}$:

\begin{align}
  C_{\ce{NaOH}}
    &=
    \frac{
      \left(\qty{\ConcHCl}{\mole\per\litre}\right)
      \left(\qty{30.40}{\milli\litre}\right)}
      {\qty{25.00}{\milli\litre}} \nonumber\\
    &=\boxed{\qty{0.1216}{\mole\per\litre}}.
  \label{eq:resultado_cnaoh}
\end{align}

As concentrações calculadas individualmente foram:

\begin{table}[H]
  \centering
  \caption{Concentrações de \ce{NaOH} calculadas para cada determinação,
  adotando $C_{\ce{HCl}}=\qty{\ConcHCl}{\mole\per\litre}$.}
  \label{tab:concentracoes}
  \begin{tabular}{
    c
    S[table-format=2.2]
    S[table-format=1.4]
  }
    \toprule
    \textbf{Determinação}
      & {\textbf{$V_{\ce{HCl}}$} (\si{\milli\litre})}
      & {\textbf{$C_{\ce{NaOH}}$} (\si{\mole\per\litre})} \\
    \midrule
    1 & 30,70 & 0,1228 \\
    2 & 30,30 & 0,1212 \\
    3 & 30,20 & 0,1208 \\
    \midrule
    \textbf{Média} & 30,40 & 0,1216 \\
    \bottomrule
  \end{tabular}
\end{table}

O fator da solução, tomando como referência a concentração nominal de
\qty{0.1}{\mole\per\litre}, foi:

\begin{align}
  f
    &=\frac{C_{\text{experimental}}}{C_{\text{nominal}}} \nonumber\\
    &=\frac{0.1216}{0.1000} \nonumber\\
    &=\boxed{1.216}.
  \label{eq:fator}
\end{align}

A diferença relativa em relação à concentração nominal foi:

\begin{align}
  \Delta C(\%)
    &=\frac{C_{\text{experimental}}-C_{\text{nominal}}}
            {C_{\text{nominal}}}\times100 \nonumber\\
    &=\frac{0.1216-0.1000}{0.1000}\times100 \nonumber\\
    &=\boxed{\qty{21.6}{\percent}}.
  \label{eq:diferenca}
\end{align}

% =========================================================
\section{Discussão}

O volume médio de \ce{HCl} consumido foi
\qty{30.40}{\milli\litre}. Considerando a concentração de
\qty{\ConcHCl}{\mole\per\litre} para o titulante, a solução de \ce{NaOH}
apresentou concentração média de \qty{0.1216}{\mole\per\litre}, valor
superior à concentração nominal de \qty{0.1}{\mole\per\litre}.

A boa concordância entre os três volumes, expressa por um desvio-padrão
relativo próximo de \qty{0.87}{\percent}, indica que a repetibilidade da
titulação foi satisfatória. Entretanto, a diferença de
\qty{21.6}{\percent} em relação ao valor nominal sugere a existência de
erro sistemático na preparação ou na concentração adotada para o
\ce{HCl}.

Entre as possíveis fontes de diferença estão a medição aproximada de
\qty{2.0}{\milli\litre} da solução concentrada de \ce{NaOH}, o uso de
material graduado para o ajuste do volume final, a concentração real da
solução concentrada de \ce{NaOH}, a absorção de água ou de
\ce{CO2} pelo reagente e eventual diferença entre a concentração nominal
e a concentração efetivamente padronizada do \ce{HCl}.

A adição de \qty{25}{\milli\litre} de água destilada no Erlenmeyer não
explica o aumento da concentração calculada. Essa água apenas diluiu
temporariamente a solução durante a titulação; a quantidade de matéria de
\ce{NaOH} permaneceu aquela contida nos
\qty{25.00}{\milli\litre} da alíquota original.

% =========================================================
\section{Conclusão}

Foi preparada uma solução de \ce{NaOH} com concentração nominal de
\qty{0.1}{\mole\per\litre} por diluição de uma solução concentrada de
aproximadamente \qty{12.4}{\mole\per\litre}. O volume calculado da solução
concentrada foi de \qty{2.016}{\milli\litre}, aproximado
experimentalmente para \qty{2.0}{\milli\litre}.

Na padronização com \ce{HCl}, os volumes gastos foram
\qty{30.70}{\milli\litre}, \qty{30.30}{\milli\litre} e
\qty{30.20}{\milli\litre}, resultando em volume médio de
\qty{30.40}{\milli\litre}. Adotando
$C_{\ce{HCl}}=\qty{\ConcHCl}{\mole\per\litre}$, a concentração média da
solução de \ce{NaOH} foi:

\begin{equation*}
  \boxed{C_{\ce{NaOH}}=\qty{0.1216}{\mole\per\litre}}.
\end{equation*}

Os \qty{25}{\milli\litre} de água adicionados ao Erlenmeyer não foram
incluídos como volume de \ce{NaOH} no cálculo, pois não acrescentaram
quantidade de matéria da base. O volume utilizado na relação
estequiométrica foi, corretamente, o da alíquota original de
\qty{25.00}{\milli\litre} de solução de \ce{NaOH}.

% =========================================================
\section*{Referência}

LAGO, Dalva Cristina Baptista do; GOIS, Jefferson Santos de; LANGONE,
Marta Antunes Pereira; SENNA, Lilian Ferreira de.
\textit{Química VIII -- Oceanografia: aulas práticas}.
Rio de Janeiro: Universidade do Estado do Rio de Janeiro,
Instituto de Química, Departamento de Química Analítica, 2017.

\end{document}
